\subsection{APPLICATION: CONVERGENCE CRITERIA OF THE SECOND KIND FOR SERIES WITH POSITIVE TERMS}

Quite often one meets series \((u_{n})\) for which \(u_{n} > 0\) from a certain point on, and \(u_{n+1}/u_{n}\) has an asymptotic expansion which is easy to determine. It is convenient, for such series, to have criteria (called criteria of the second kind) allowing one to determine whether the series is convergent from the form of \(u_{n+1}/u_{n}\) alone. The following is such a criterion:

PROPOSITION 7 (``Raabe's test''). Let \((u_{n})\) be a series with terms \textgreater{} 0 from some point on. If, from a certain stage, \(u_{n+1}/u_{n} \leqslant 1 - \frac{\alpha}{n}\) for some \(\alpha > 1\) , then the series \((u_{n})\) converges; if, from a certain point on, \(u_{n+1}/u_{n} \geqslant 1 - \frac{1}{n}\) , then the series \((u_{n})\) has infinite sum.

Indeed, if \(u_{n+1} / u_n \leqslant 1 - \frac{\alpha}{n}\) with \(\alpha > 1\) , for all \(n \geqslant n_0\) , then one has \(u_n \preccurlyeq p_n = \prod_{k=n_0}^{n} \left(1 - \frac{\alpha}{k}\right)\) . Now, \(\log \left(1 - \frac{\alpha}{n}\right) = -\frac{\alpha}{n} - \frac{\alpha^2}{2n^2} + o\left(\frac{1}{n^2}\right)\) , whence \(\log p_n =\) \(-\alpha \log n + k + o(1/n) (k \text{ constant}), \text{ and } p_n \sim e^k \frac{1}{n^\alpha}; \text{ since } \alpha > 1 \text{ the logarithmic criterion of order 0 allows one to finish.}\)

If, on the other hand, \(u_{n + 1} / u_n \geqslant 1 - \frac{1}{n}\) from a certain point, then the same calculation proves that \(u_n \succcurlyeq \frac{1}{n}\) , whence the proposition.

One can prove in the same way, using the logarithmic criteria of order \textgreater{} 0, the following criterion of the second kind:

PROPOSITION 8. Let \((u_n)\) be a series with terms \(>0\) from a certain point on. If, from a certain point, one has

\[
\frac {u _ {n + 1}}{u _ {n}} \leqslant 1 - \frac {1}{n} - \frac {1}{n l _ {1} (n)} - \dots - \frac {1}{n l _ {1} (n) l _ {2} (n) \dots l _ {p - 1} (n)} - \frac {\alpha}{n l _ {1} (n) l _ {2} (n) \dots l _ {p} (n)}
\]

for some \(\alpha > 1\) , then the series \((u_n)\) converges; if, from some point on, one has

\[
\frac {u _ {n + 1}}{u _ {n}} \geqslant 1 - \frac {1}{n} - \frac {1}{n l _ {1} (n)} - \dots - \frac {1}{n l _ {1} (n) l _ {2} (n) \dots l _ {p} (n)}
\]

then the series \((u_{n})\) has an infinite sum.

Example. Consider the hypergeometric series, with general term

\[
u _ {n} = \frac {\alpha (\alpha + 1) \dots (\alpha + n - 1) \beta (\beta + 1) \dots (\beta + n - 1)}{1 . 2 \dots n . \gamma (\gamma + 1) \dots (\gamma + n - 1)}
\]

where \(\alpha, \beta, \gamma\) are arbitrary real numbers, not integers \(\leqslant 0\) ; it is clear that \(u_{n}\) is \(>0\) from a certain stage, or \(< 0\) from a certain stage on. One has

\[
\begin{array}{l} \frac {u _ {n + 1}}{u _ {n}} = \frac {(\alpha + n) (\beta + n)}{(n + 1) (\gamma + n)} \\ \qquad = \left(1 + \frac {\alpha + \beta}{n} + \frac {\alpha \beta}{n ^ {2}}\right) \left(1 + \frac {\gamma + 1}{n} + \frac {\gamma}{n ^ {2}}\right) ^ {- 1} \\ \qquad = 1 + \frac {\alpha + \beta - \gamma - 1}{n} \\ \qquad + \frac {\alpha \beta - (\alpha + \beta) (\gamma + 1) + \gamma^ {2} + \gamma + 1}{n ^ {2}} + o \left(\frac {1}{n ^ {2}}\right). \end{array}
\]

Raabe's test shows that the series converges for \(\alpha +\beta <  \gamma\) , and has infinite sum for \(\alpha +\beta >\gamma\) ; when \(\alpha +\beta = \gamma\) the series also has an infinite sum, as is shown by prop. 8.

Remarks. 1) As a particular instance of Raabe's test one sees that if \(\limsup_{n\to \infty}u_{n + 1} / u_n\) \(< 1\) the series \((u_{n})\) converges; if, on the other hand, \(\liminf_{n\to \infty}u_{n + 1} / u_n > 1\) , the series \((u_{n})\) has an infinite sum (d'Alembert's test).

\begin{enumerate}
\def\labelenumi{\arabic{enumi})}
\setcounter{enumi}{1}
\tightlist
\item
  The criteria of the second kind can be applied only to series whose general term behaves in a very regular way as n tends to \(+\infty\) ; in other words, their scope is much more restricted than those of the logarithmic criteria, and it would be a blunder to try to use them beyond the special cases to which they are particularly suited. For example, for the series \((u_{n})\) defined by \(u_{2m} = 2^{-m}\) , \(u_{2m + 1} = 3^{-m}\) one has \(u_{2m + 1} / u_{2m} = \left(\frac{2}{3}\right)^{m}\) , \(u_{2m + 2} / u_{2m + 1} = \frac{1}{2}\left(\frac{3}{2}\right)^{m}\) ; the first of these ratios tends to 0 and the second to \(+\infty\) as \(m\) increases indefinitely, so no criterion of the second kind is applicable; nevertheless, since \(u_{n} \preccurlyeq 2^{-n / 2}\) , it is immediate that the series converges.
\end{enumerate}

Even when \(u_{n+1} / u_n\) has a simple expression, a direct evaluation of a principal part for \(u_n\) often leads to a result as quickly as the criteria of the second kind. For example, for the hypergeometric series, Stirling's formula shows immediately that \(u_n \sim a n^{\alpha + \beta - \gamma - 1}\) , where \(a\) is a constant \(\neq 0\) , and the logarithmic criterion of order 0 is then applicable.

\appendix
\renewcommand{\thechapter}{}
\renewcommand{\thesection}{\arabic{section}}
